\documentclass[10.5pt,a4paper]{article}

\usepackage{fontspec}
\usepackage[a4paper,margin=2.2cm,headheight=15pt]{geometry}
\usepackage{booktabs}
\usepackage{tabularx}
\usepackage{array}
\usepackage{float}
\usepackage{microtype}
\usepackage{xcolor}
\usepackage{enumitem}
\usepackage{titlesec}
\usepackage{fancyhdr}
\usepackage{xurl}
\usepackage{amsmath}
\usepackage{hyperref}

\definecolor{navy}{HTML}{17324D}
\definecolor{blue}{HTML}{1F6FEB}
\definecolor{pale}{HTML}{F4F7FA}
\hypersetup{colorlinks=true,linkcolor=blue,urlcolor=blue,citecolor=blue}
\setlength{\parskip}{0.28em}
\setlength{\parindent}{0pt}
\setlist{nosep,leftmargin=1.5em}
\titlespacing*{\section}{0pt}{1.4ex plus .2ex}{0.6ex}
\titlespacing*{\subsection}{0pt}{1.0ex plus .2ex}{0.4ex}
\renewcommand{\arraystretch}{1.18}
\raggedbottom
% Generated by scripts/render_report_data.py. Do not edit by hand.
\newcommand{\BestModel}{Ensemble}
\newcommand{\BestMAPE}{1.855\%}
\newcommand{\BestMASE}{0.442}
\newcommand{\BestMAE}{997.7}
\newcommand{\FinalRole}{retrospective\_final}
\newcommand{\FinalRoleLabel}{Retrospective; previously inspected}
\newcommand{\FinalStart}{2026-01-01}
\newcommand{\FinalEnd}{2026-08-04}
\newcommand{\FinalN}{216}
\newcommand{\SourceRevision}{e166beb21e3bc332989a73a0ee002f1fc5d05b70}
\newcommand{\BundleFingerprint}{c5886267e46429d2e47ee218917d62c4037a49c6906dba6e7ea02c9a9886483e}
\newcommand{\DailyProtocolFingerprint}{fed932e555077113ad2b8bba15935e0d02b4323dd2a6c8ece97e4d475c09bcb7}
\newcommand{\HourlyProtocolFingerprint}{4ece152109b0d6e895aca13b2a4895fc799a564f6f307c36f184860854e0bfa1}
\newcommand{\EnsembleMAPE}{1.855\%}
\newcommand{\EnsembleMASE}{0.442}
\newcommand{\EnsembleMAE}{997.7}
\newcommand{\HourlySelectedModel}{Residual hybrid}
\newcommand{\DailyAnchorModel}{LightGBM}
\newcommand{\HourlyValidationSelectedMAE}{1463.6}
\newcommand{\HourlyValidationLightGBMMAE}{1466.6}
\newcommand{\HourlyTestSelectedMAE}{1562.1}
\newcommand{\HourlyTestLightGBMMAE}{1558.0}
\newcommand{\BootstrapDifference}{1.6}
\newcommand{\BootstrapLower}{-44.5}
\newcommand{\BootstrapUpper}{46.7}
\newcommand{\BootstrapProbability}{47.74\%}
\newcommand{\BootstrapN}{216}
\newcommand{\RawHourlyMAE}{1624.9}
\newcommand{\ReconciledHourlyMAE}{1562.1}
\newcommand{\HourlyN}{5183}
\newcommand{\ReconciliationGain}{3.86\%}
\newcommand{\SymmetricNominalLevel}{90\%}
\newcommand{\SymmetricCoverage}{85.18\%}
\newcommand{\SymmetricWidth}{5612.0}
\newcommand{\SymmetricIntervalScore}{8837.3}
\newcommand{\SymmetricN}{5183}
\newcommand{\AdaptiveNominalLevel}{90\%}
\newcommand{\AdaptiveCoverage}{89.95\%}
\newcommand{\AdaptiveWidth}{6545.0}
\newcommand{\AdaptiveIntervalScore}{8273.5}
\newcommand{\AdaptiveN}{5183}
\newcommand{\CQRNominalLevel}{90\%}
\newcommand{\CQRCoverage}{86.67\%}
\newcommand{\CQRWidth}{7556.4}
\newcommand{\CQRIntervalScore}{10519.8}
\newcommand{\CQRN}{5183}
\newcommand{\EmpiricalRetrospectiveCoverageLabel}{Empirical retrospective coverage.}
\newcommand{\PrequentialMonitoringCoverageLabel}{Prequential monitoring evidence; no unconditional time-series coverage guarantee.}
\newcommand{\ArtifactReplayError}{7.28e-12}

\newcommand{\ValidationRows}{%
SARIMAX & 1640.2 & 1182.1 & 2.276 & 0.532 \\
XGBoost & 1421.6 & 1077.9 & 2.081 & 0.486 \\
LightGBM & 1463.2 & 1096.1 & 2.113 & 0.494 \\
Random Forest & 1779.7 & 1191.0 & 2.308 & 0.536 \\
Naive (t-1) & 5189.4 & 3682.2 & 7.190 & 1.659 \\
Seasonal naive (t-7) & 3832.2 & 2519.6 & 4.873 & 1.135 \\
Ensemble & 1394.9 & 1030.4 & 1.984 & 0.464 \\
}

\newcommand{\CalibrationRows}{%
SARIMAX & 1308.5 & 1000.6 & 1.906 & 0.437 \\
XGBoost & 1202.9 & 943.6 & 1.807 & 0.412 \\
LightGBM & 1244.5 & 971.2 & 1.860 & 0.424 \\
Random Forest & 1377.8 & 1012.3 & 1.930 & 0.442 \\
Naive (t-1) & 4970.2 & 3471.1 & 6.806 & 1.516 \\
Seasonal naive (t-7) & 2721.1 & 1813.9 & 3.433 & 0.792 \\
Ensemble & 1149.7 & 918.3 & 1.752 & 0.401 \\
}

\newcommand{\TestRows}{%
SARIMAX & 1465.6 & 1090.6 & 2.040 & 0.484 \\
XGBoost & 1371.6 & 1074.2 & 1.999 & 0.476 \\
LightGBM & 1413.8 & 1127.8 & 2.095 & 0.500 \\
Random Forest & 1567.4 & 1154.0 & 2.139 & 0.512 \\
Naive (t-1) & 5160.6 & 3658.1 & 7.025 & 1.622 \\
Seasonal naive (t-7) & 4189.9 & 2978.8 & 5.575 & 1.321 \\
Ensemble & 1302.8 & 997.7 & 1.855 & 0.442 \\
}

\newcommand{\RollingRows}{%
2024-H1 & 2.412 & 2.103 & 2.236 & 2.471 & 5.314 \\
2024-H2 & 1.995 & 1.796 & 1.760 & 2.114 & 3.980 \\
2025-H1 & 2.228 & 1.910 & 1.905 & 2.160 & 5.336 \\
2025-H2 & 1.906 & 1.807 & 1.860 & 1.930 & 3.433 \\
}

\newcommand{\IntervalRows}{%
SARIMAX & Fixed & 95\% & 0.95 & 0.898 & 4734.4 & 8496.6 & 216 \\
SARIMAX & Adaptive & 95\% & 0.95 & 0.958 & 6053.9 & 8082.5 & 216 \\
XGBoost & Fixed & 95\% & 0.95 & 0.912 & 4571.8 & 7288.9 & 216 \\
XGBoost & Adaptive & 95\% & 0.95 & 0.958 & 5608.5 & 7100.4 & 216 \\
LightGBM & Fixed & 95\% & 0.95 & 0.931 & 5275.1 & 6849.0 & 216 \\
LightGBM & Adaptive & 95\% & 0.95 & 0.958 & 5641.7 & 6827.8 & 216 \\
Random Forest & Fixed & 95\% & 0.95 & 0.940 & 5611.9 & 9211.1 & 216 \\
Random Forest & Adaptive & 95\% & 0.95 & 0.963 & 7229.9 & 8906.4 & 216 \\
Naive (t-1) & Fixed & 95\% & 0.95 & 0.958 & 21680.3 & 22625.7 & 216 \\
Naive (t-1) & Adaptive & 95\% & 0.95 & 0.949 & 21927.3 & 23034.7 & 216 \\
Seasonal naive (t-7) & Fixed & 95\% & 0.95 & 0.870 & 10173.0 & 30556.2 & 216 \\
Seasonal naive (t-7) & Adaptive & 95\% & 0.95 & 0.954 & 20461.1 & 26972.7 & 216 \\
Ensemble & Fixed & 95\% & 0.95 & 0.926 & 4600.3 & 6647.4 & 216 \\
Ensemble & Adaptive & 95\% & 0.95 & 0.963 & 4959.7 & 6624.0 & 216 \\
}

\newcommand{\WeatherRows}{%
SARIMAX & 2.040 & 2.058 & -0.018 \\
XGBoost & 1.999 & 2.030 & -0.031 \\
LightGBM & 2.095 & 2.126 & -0.031 \\
Random Forest & 2.139 & 2.230 & -0.091 \\
Ensemble & 1.855 & 1.900 & -0.045 \\
}

\newcommand{\GeneralizationRows}{%
SARIMAX & 2.276 & 2.040 & -0.236 \\
XGBoost & 2.081 & 1.999 & -0.082 \\
LightGBM & 2.113 & 2.095 & -0.017 \\
Random Forest & 2.308 & 2.139 & -0.169 \\
Naive (t-1) & 7.190 & 7.025 & -0.164 \\
Seasonal naive (t-7) & 4.873 & 5.575 & +0.703 \\
Ensemble & 1.984 & 1.855 & -0.128 \\
}


\pagestyle{fancy}
\fancyhf{}
\lhead{\textcolor{navy}{German Grid-Load Forecasting}}
\rhead{\textcolor{navy}{Technical Report}}
\cfoot{\thepage}
\renewcommand{\headrulewidth}{0.3pt}

\newcommand{\callout}[2]{%
  \begin{center}
  \colorbox{pale}{\parbox{0.92\textwidth}{\textbf{\textcolor{navy}{#1}}\quad #2}}
  \end{center}
}
\newcommand{\metric}[2]{%
  \begin{minipage}[t]{0.23\textwidth}
  \centering
  {\Large\bfseries\textcolor{blue}{#1}}\par
  \small #2
  \end{minipage}
}

\begin{document}

\begin{titlepage}
\color{navy}
\vspace*{1.2cm}
{\Large\textbf{TECHNICAL REPORT}}\par
\vspace{0.6cm}
{\fontsize{30}{36}\selectfont\bfseries German Grid-Load Forecasting\par}
\vspace{0.5cm}
{\Large Daily and hourly day-ahead forecasts}\par
\vspace{1.2cm}
\colorbox{pale}{\parbox{0.88\textwidth}{
The system combines a daily forecast and an hourly profile. It uses
hierarchical reconciliation and conformal prediction intervals. MLflow stores
the complete reconciled model.
}}
\vfill
\begin{tabular}{@{}ll}
\textbf{Author} & chiki · \href{mailto:work@chiki.dev}{work@chiki.dev}\\
\end{tabular}
\vspace{0.8cm}
\end{titlepage}

\section*{Summary}

The system forecasts realized German electricity grid load. It produces daily
and hourly day-ahead forecasts.

% loadfc:generated-start
Release results are generated from the checked source revision.
% loadfc:generated-end

\section{Data and evaluation}

SMARD supplies realized grid-load data. Open-Meteo supplies archived weather
forecasts.

The target is realized electricity consumption in MW. Residual load is outside
the report scope.

Weather features use forecasts from 24 hours before valid time. Earlier data
uses previous-day observed weather.

The configured chronological stages are training, validation, calibration, and
retrospective final. Validation selects models before later evidence is inspected.

\section{Forecast architecture}

\begin{center}
\setlength{\fboxsep}{8pt}
\begin{tabular}{c@{\hspace{0.35em}$\longrightarrow$\hspace{0.35em}}c@{\hspace{0.35em}$\longrightarrow$\hspace{0.35em}}c}
\fbox{\parbox{0.24\textwidth}{\centering\textbf{Inputs}\\load\\weather}}
&
\fbox{\parbox{0.24\textwidth}{\centering\textbf{Models}\\hourly profile\\daily mean}}
&
\fbox{\parbox{0.24\textwidth}{\centering\textbf{Outputs}\\reconciliation\\intervals · MLflow}}
\end{tabular}
\end{center}

The hourly evaluation compares residual hybrid, direct LightGBM, and direct
Ridge. Each model predicts valid hours directly.

The features include horizon, lag-24, lag-168, calendar, weather, and Fourier
terms. Direct prediction prevents recursive error propagation.

Validation MAE selects the daily anchor. The final model uses only data before
calibration.

\section{Reconciliation}

For local day $d$ and hour $h$, the system uses:

\[
\hat{y}^{rec}_{d,h}
=
\hat{y}^{hourly}_{d,h}
\times
\frac{\hat{y}^{daily}_{d}}
{\frac{1}{|H_d|}\sum_{j\in H_d}\hat{y}^{hourly}_{d,j}}.
\]

The operation keeps the relative hourly shape. The reconciled mean equals the
daily forecast.

The code supports complete local days with 23, 24, or 25 hours. It rejects
invalid predictions and incomplete days.

The generated release evidence reports raw and reconciled MAE from the same
retrospective-final forecast rows.

\section{Model selection}

The release keeps the validation-selected model. Later candidate comparisons
are retrospective evidence and do not change that selection.

The paired bootstrap uses one MAE difference for each local day and resamples
the paired day blocks. The generated release evidence reports its result.

\section{Prediction intervals}

The point, daily, and quantile models use data before calibration. The models
stay fixed during calibration and retrospective-final evaluation.

Reconciliation occurs before score calculation. Both interval methods use the
same reconciled point forecasts.

The current implementation keeps all models frozen. Coverage is empirical; adaptive intervals are prequential monitoring without an unconditional time-series guarantee.

\section{Controls}

\subsection{MLflow}

\texttt{ReconciledForecaster} contains both fitted models and both feature
schemas. It also contains the reconciliation operation.

Artifact replay is checked against the persisted evaluation predictions.

\subsection{Data validation}

\begin{tabularx}{\textwidth}{>{\bfseries}l X}
Load & Values must be finite and positive.\\
Weather & Values must be within configured physical ranges.\\
Time & UTC timestamps must be unique, complete, and on the hourly grid.\\
Telemetry & Prediction and interval indexes must match before the join.\\
\end{tabularx}

\subsection{Datetime resolution}

Raw \texttt{.asi8} values depend on the pandas storage unit. The feature code
uses \texttt{Timedelta(hours=1)} to calculate elapsed hours.

A regression test uses \texttt{datetime64[us]}. It verifies the time scale and
the weekly Fourier phase.

\section{Verification and limits}

\begin{verbatim}
uv run ruff check src scripts tests
uv run mypy src
uv run pytest --cov=loadfc --cov-report=term-missing --cov-fail-under=80
uv run python scripts/validate_results.py
\end{verbatim}

Five cities represent national weather exposure. Distribution shift can change
conformal coverage.

A production service requires scheduling, access control, service-level
objectives, and rollback rules.

\section*{References}

\begin{itemize}
\footnotesize
  \item Bundesnetzagentur, SMARD: \url{https://www.smard.de/en}
  \item Open-Meteo Previous Runs API:
        \url{https://open-meteo.com/en/docs/previous-runs-api}
  \item Hyndman et al., ``Optimal Combination Forecasts for Hierarchical Time Series''.
  \item Romano et al., ``Conformalized Quantile Regression''.
  \item Gibbs and Candès, ``Adaptive Conformal Inference Under Distribution Shift''.
  \item Ke et al., ``LightGBM: A Highly Efficient Gradient Boosting Decision Tree''.
\end{itemize}

\end{document}
